\paragraph{事件数字的 Gauss--Kuzmin 普适律与 Lee--Yang 边界的测度桥接}
在 \(R,L\) 的事件记录口径下，事件数字序列 \(\mathbf a=(a_1,a_2,\dots)\) 可视为“分支扩张 + 取整”的离散影子；
其经验分布若满足 Gauss 测度的数字律，则可把事件统计、模曲面动力学与证据配分函数的复奇点边界置于同一测度框架之中。

\begin{conjecture}[事件数字的 Gauss--Kuzmin 极限与模曲面等分布]\label{conj:xi-gauss-kuzmin-event-digit-bridge}
沿用结论链中出现的生成元
\(
R=\bigl(\begin{smallmatrix}1&1\\0&1\end{smallmatrix}\bigr)
\)、
\(
L=\bigl(\begin{smallmatrix}1&0\\1&1\end{smallmatrix}\bigr)
\)，
并对有限前缀 \(\mathbf a_{1:T}=(a_1,\dots,a_T)\in(\NN_{\ge 1})^T\) 定义
\[
M_T:=M(\mathbf a_{1:T})=R^{a_1}L\cdots R^{a_T}L\in \mathrm{SL}_2(\ZZ_{>0}),
\qquad
\tau_T:=M_T\cdot i\in\mathbb H.
\]
设事件数字 \((a_t)_{t\ge 1}\) 由某一“无结构基线”微观驱动诱导（例如独立同分布输入经折叠/读出机制推前的事件记录）。
则猜想：
\begin{enumerate}
  \item \textbf{Gauss--Kuzmin 数字律：}
  对每个 \(k\ge 1\)，经验频率满足
  \[
  \frac{1}{T}\#\{1\le t\le T:\ a_t=k\}\ \longrightarrow\ 
  \PP_{\mathrm{GK}}(k)
  :=\frac{1}{\log 2}\,\log\!\left(1+\frac{1}{k(k+2)}\right),
  \]
  即与经典 Gauss 测度下连分数数字分布一致。
  \item \textbf{模曲面等分布：}
  \((\tau_T)_{T\ge 1}\) 在 \(\mathrm{PSL}_2(\ZZ)\backslash\mathbb H\) 上按 Gauss 测度诱导的不变概率测度等分布。
  \item \textbf{Lee--Yang 边界的测度锁定：}
  若证据配分函数 \(Z_n(z)\) 的 Lee--Yang 极限边界 \(\mathcal K\subset\CC^\times\) 存在（猜想 \ref{conj:xi-leyang-boundary-capacity-coherence}），则 \(\mathcal K\) 的势函数/容量常数可由上述 Gauss 极限测度下的 Lyapunov 自由能（定理 \ref{thm:xi-leyang-evidence-ising-lyapunov}）决定；
  因而事件数字律、模点轨道与 Lee--Yang 边界在测度层面闭合为同一对象的三个投影。
\end{enumerate}
\end{conjecture}

